\subsection{Haar probability and continuous laws}\label{subsec:measures-preliminaries-laws}
\paraid{p-c03d5071df48}
We use Haar probability to average over the isometry group.

\paraid{p-d2f74a62001a}
By the next lemma,
the full isometry group is compact,
so we can apply
the Haar probability theorem.

\begin{lemma}\label{lem:compact-isometry-group}
\paraid{p-bd8a5b8d04a8}
Let
$\MetricSpace{\BaseCarrier}{\MetricSymbol}$
be a nonempty compact metric space.
Then
$\Isom\MetricSpace{\BaseCarrier}{\MetricSymbol}$,
equipped with the uniform topology,
is a compact metrizable topological group.
\end{lemma}

\begin{proof}
\paraid{p-e174430bf1e9}
Let
$\IsometryGroup=\Isom\MetricSpace{\BaseCarrier}{\MetricSymbol}$.
For
$\Mapping,\SecondMapping\in\ContinuousFunctions(\BaseCarrier,\BaseCarrier)$,
define the uniform metric by
\[
 \UniformMetric(\Mapping,\SecondMapping)
 =\sup_{\BasePoint\in\BaseCarrier}
    \MetricSymbol(\Mapping(\BasePoint),\SecondMapping(\BasePoint)).
\]
Since every element of
$\IsometryGroup$
is an isometry and
$\BaseCarrier$
is compact,
the Arzel\`a--Ascoli theorem (\Cref{thm:arzela-ascoli}) implies that
$\IsometryGroup$
has compact closure in
$\ContinuousFunctions(\BaseCarrier,\BaseCarrier)$.

\paraid{p-6888b18e99f0}
A uniform limit
$\Mapping\in\ContinuousFunctions(\BaseCarrier,\BaseCarrier)$
of elements of
$\IsometryGroup$
preserves distances and is surjective by
\cite[Theorem~1.6.14]{BuragoBuragoIvanov2001}.
Thus
$\IsometryGroup$
is closed and compact.
Let
$\GroupElement_\SequenceIndex\to\GroupElement$
and
$\SecondGroupElement_\SequenceIndex\to\SecondGroupElement$
be two convergent sequences in
$\IsometryGroup$.
From the estimates
\[
 \UniformMetric(\GroupElement_\SequenceIndex\circ\SecondGroupElement_\SequenceIndex,
               \GroupElement\circ\SecondGroupElement)
 \leq\UniformMetric(\GroupElement_\SequenceIndex,\GroupElement)
       +\UniformMetric(\SecondGroupElement_\SequenceIndex,\SecondGroupElement),
\]
and
\[
 \UniformMetric(\GroupElement_\SequenceIndex^{-1},\GroupElement^{-1})
 =\UniformMetric(\GroupElement_\SequenceIndex,\GroupElement),
\]
it follows  that composition and inversion are continuous.
Hence
$\IsometryGroup$
is a compact metrizable topological group.
\end{proof}

\begin{theorem}[{Haar probability, \cite[Theorem~5.14 and Chapter~5, Section~3]{DiestelSpalsbury2014}}]\label{thm:haar}
\paraid{p-c37dcbf3e97d}
Every compact metrizable topological group
$\ActingGroup$
admits a unique left invariant Radon probability
$\HaarProbability\in\ProbabilityMeasures(\ActingGroup)$.
Thus
\[
 \HaarProbability(\ActingGroup)=1,
\]
and for every
$\GroupElement\in\ActingGroup$
and every
$\BorelSubset\in\BorelSets(\ActingGroup)$,
we have
\[
 \HaarProbability(\GroupElement\BorelSubset)=\HaarProbability(\BorelSubset).
\]
This probability is also right invariant and has full support.
\end{theorem}

\paraid{p-09eb6dcfaa8e}
We apply the following consequence of Valov's theorem
\cite[arXiv v2, Theorem~1.1]{Valov2009}.
The probabilities in this formulation are Radon probabilities.
The theorem imposes no compact-support restriction.
For Polish spaces, the same conclusion follows from the earlier theorem
of Banakh and Radul \cite[Theorem~1.2, p.~20]{BanakhRadul1999}.

\begin{theorem}[{Valov, \cite[arXiv v2, Theorem~1.1]{Valov2009}}]\label{thm:valov}
\paraid{p-ea9930a1ae3f}
Let
$\OpenSurjection \colon \LawTotalSpace \to \LawBaseSpace $
be an open continuous surjection between completely metrizable spaces.
Then there is a continuous map
$\LawSection\colon \LawBaseSpace \to\ProbabilityMeasures(\LawTotalSpace )$
such that
$\OpenSurjection \Pushforward\LawSection(\LawBasePoint )=\DiracProbability _\LawBasePoint $
for every
$\LawBasePoint \in \LawBaseSpace $.
\end{theorem}

\paraid{p-07d85ea78312}
Here
$\DiracProbability _\LawBasePoint $
is the Dirac probability at
$\LawBasePoint $.
Valov's theorem yields a continuous right inverse of the induced
map on probabilities.
Composing this right inverse with
$\LawBasePoint \mapsto\DiracProbability _\LawBasePoint $
proves this formulation.

\begin{remark}\label{rem:laws-averaging-operator}
\paraid{p-2db988c38721}
The law in \Cref{thm:valov} also defines a linear averaging operator.
Assume that
$\LawBaseSpace\ne\emptyset$,
and define the averaging operator
\[
 \AveragingOperator\colon
 \BoundedContinuousFunctions(\LawTotalSpace,\RealNumbers)
 \to\BoundedContinuousFunctions(\LawBaseSpace,\RealNumbers).
\]
For every
$\TestFunction\in\BoundedContinuousFunctions(\LawTotalSpace,\RealNumbers)$
and
$\LawBasePoint\in\LawBaseSpace$,
put
\[
 (\AveragingOperator\TestFunction)(\LawBasePoint)
 =\int_{\LawTotalSpace}\TestFunction(\FiberPoint)
   \,\IntegrationDifferential\LawSection(\LawBasePoint)(\FiberPoint).
\]
Since
$\LawSection$
is a weakly continuous family of probabilities,
$\AveragingOperator$
is a positive linear operator satisfying
\begin{equation}\label{eq:law-averaging-identities-1}
 \AveragingOperator1=1.
\end{equation}
For every
$\TestFunction\in\BoundedContinuousFunctions(\LawTotalSpace,\RealNumbers)$,
\begin{equation}\label{eq:law-averaging-identities-2}
 \|\AveragingOperator\TestFunction\|_\infty\leq\|\TestFunction\|_\infty.
\end{equation}
The fiber support condition implies that,
for every
$f\in\BoundedContinuousFunctions(\LawBaseSpace,\RealNumbers)$,
\begin{equation}\label{eq:law-averaging-identities-3}
 \AveragingOperator(f\circ\OpenSurjection)=f.
\end{equation}
Define
$\OpenSurjection^*\colon
 \BoundedContinuousFunctions(\LawBaseSpace,\RealNumbers)
 \to\BoundedContinuousFunctions(\LawTotalSpace,\RealNumbers)$
by
$\OpenSurjection^*f=f\circ\OpenSurjection$.
By \eqref{eq:law-averaging-identities-1}--\eqref{eq:law-averaging-identities-3},
$\OpenSurjection^*\AveragingOperator$
is a projection of norm one onto
$\OpenSurjection^*\BoundedContinuousFunctions(\LawBaseSpace,\RealNumbers)$.
\end{remark}
